%%
%% This is file `template-8s.tex',
%% generated with the docstrip utility.
%%
%% The original source files were:
%%
%% template.raw  (with options: `8s')
%% 
%% Template for the LaTeX class aipproc.
%% 
%% (C) 1998,2000,2001 American Institute of Physics and Frank Mittelbach
%% All rights reserved
%% 
%%
%% $Id: template.raw,v 1.11 2004/10/31 08:06:14 frank Exp $
%%

%%%%%%%%%%%%%%%%%%%%%%%%%%%%%%%%%%%%%%%%%%%%
%% Please remove the next line of code if you
%% are satisfied that your installation is
%% complete and working.
%%
%% It is only there to help you in detecting
%% potential problems.
%%%%%%%%%%%%%%%%%%%%%%%%%%%%%%%%%%%%%%%%%%%%

%\input{aipcheck}

%%%%%%%%%%%%%%%%%%%%%%%%%%%%%%%%%%%%%%%%%%%%
%% SELECT THE LAYOUT
%%
%% The class supports further options.
%% See aipguide.pdf for details.
%%
%%%%%%%%%%%%%%%%%%%%%%%%%%%%%%%%%%%%%%%%%%%%

\documentclass[
   ,final            % use final for the camera ready runs
%% ,draft            % use draft while you are working on the paper
%%  ,numberedheadings % uncomment this option for numbered sections
%%  ,                 % add further options here if necessary
  ]
  {aipproc}

\layoutstyle{8x11single}
\usepackage{palatino,url,amsmath,amssymb,enumerate,inputenc}
%\usepackage[latin1]{inputenc} % From LaTeX distribution
\usepackage{calc}         % From LaTeX distribution
\usepackage{graphicx}     % From LaTeX distribution
\usepackage{ifthen}       % From LaTeX distribution
\usepackage{subfigure}    % From CTAN/macros/latex/contrib/supported/subfigure
\usepackage{pst-all}      % From PSTricks
\usepackage{pst-poly}     % From pstricks/contrib/pst-poly
\usepackage{multido}      % From PSTricks
%%%%%%%%%%%%%%%%%%%%%%%%%%%%%%%%%%%%%%%%%%%%
%% FRONTMATTER
%%%%%%%%%%%%%%%%%%%%%%%%%%%%%%%%%%%%%%%%%%%%

\begin{document}

\title[Quantum Systems \& Self-referentiality]{Why do Quantum Systems Implement Self-referential
Logic? A Simple Question With a Catastrophic Answer.}

\classification{12.10.-g,04.20,-q,03.67.Lx,03.65.Ta}
\keywords      {Causality protection, self-referentiality, entanglement, spacetime dimensionality,Standard Model}

\author{John Small}{
  address={40 West Street, Faversham, Kent, ME13 7JG, UK \\
  email:jds@mindoro-marine.co.uk}
}

\begin{abstract}
Simple self-referential algorithms can model the essential logic of quantum systems and the quantum formalism is sufficiently rich to be able to model %%@
the output states of
self-referential alogrithms. We identify a key feature of self-referentiality and show that both
quantum mechanics and general relativity exhibit this feature in distinct yet complementary ways. This implies
that self-referentiality must be a foundational principle, but causality protection prohibits it being directly
observed. Examining two facets of causality protection we find that quantum mechanics implements one and
general relativity the other. The implication is that there can be no mathematically consistent theory for quantum
gravity. However we can gain understanding of the constraints imposed on physical law when viewing a noncausal
reality as if it were a causal sequence. We derive the 3+1 dimensionality of spacetime and
quantum non-locality and also possible rationale for the $U(1) \times SU(2) \times SU(3)$ symmetry of the
Standard Model. We also identify a number of key open problems. 
\end{abstract}

\maketitle

%%%%%%%%%%%%%%%%%%%%%%%%%%%%%%%%%%%%%%%%%%%%
%% MAINMATTER
%%%%%%%%%%%%%%%%%%%%%%%%%%%%%%%%%%%%%%%%%%%%

\section{1 INTRODUCTION }
\par
It is widely held that the task of finding a generally accepted Unified Field Theory, or even the more
limited task of finding a theory of quantum gravity, is presently held up by lack of some key insight into
the nature of both quantum theory and general relativity. There are key aspects of quantum theory which
disagree with key aspects of general relativity. Some authors contend that general relativity must yield to
quantum theory, and others that quantum theory must yield to general relativity, still others contend that
both must be derivative of some underlying theory, but we cannot yet specify what that is ( see \citet*{phys_meets_philosophy} and \citet{Isham_92} %%@
for a good introduction). Both
these theories tell us something important about the true nature of the physical world, but but seem to
express that in different and irreconcilable ways. It is sometimes thought that a theory of quantum gravity, or a Unified Field Theory
is a case of finding the right mathematical formula. I believe the task, which this paper takes up, is
to find the right concepts to which we can apply our mathematical tools. As John Baez says in \citep[chap.~8]{phys_meets_philosophy}
\begin{quotation}General relativity and quantum field theory are based on some profound insights about the
nature of reality. These insights are crystallized in the form of mathematics, but there is a limit
to how much progress we can make by just playing around with this mathematics. We need to
go back to the insights behind general relativity and quantum field theory, learn to hold them
together in our minds, and dare to imagine a world more strange, more beautiful, but ultimately
more \emph{reasonable} than our current theories of it. For this daunting task, philosophical reflection is
bound to be of help
\end{quotation}
\par
The purpose of this paper is to make an attempt to find those insights to which we can apply mathematics. 
Therefore the word to equation ratio is very much in favour of the former. We look directly at the most paradoxical features of these theories to %%@
learn the deepest lessons. Analysis of these key insights allows us to make a connection with mathematics which does then lead to useful
predictions about the physical world although the implication is that there can be no Unified Field Theory.  The two main sections of the paper are %%@
(2) and (3), the first of these collects
evidence from current theory that an element of non-computability in both quantum mechanics and general relativity
but appearing in different guises.  From this we deduce that such non-computability must be a fundamental
feature of Nature and hence a mathematically consistent Unified Field Theory is not possible. In section (3) we look at the fact that 
we do not directly find evidence for such non-computablity because physical law appears to be constructed
so as to disallow the violation of causality which would ensue. I propose to turn this around and ask "If
we explicitly disallow non-computability what what form must physical law take?". This part of
the paper examines those constraints which must be imposed on our view of such a non-computable, and
hence non-causal reality if we to try to represent it as a strictly causal sequence of computable steps. In this
case making use of computational models of reality is useful because it enables us to more easily deal with
the notion of non-computability. From this we can gain insight into problems which are otherwise hard to
understand and we can start to make some useful deductions about the form of physical law. However
we are not seeking a full formal theory at present, at this level we are merely seeking out possible plausible
arguments that might point us in the right direction to find a full theory, absolute rigor will have to wait.
\section{2 LESSONS FROM CURRENT THEORY }
\subsection{2.1 The lesson from quantum mechanics }
\par
The recent advance in quantum computation has very much improved our understanding of the logic
of quantum mechanics. The present literature on quantum computation is too vast to review here; however a foundational
paper by \citet{Deutsch_91} discusses the use of quantum logic in cases where closed time-like loops are present
and finds that the quantum formalism can model the states encountered in such systems. This seems to be a very important result that has not received %%@
the attention it deserves and so before progressing to the main body of the argument we need to look at the points he brought out in that paper. In %%@
summary we envisage a situation where, by some means, we can create a closed time-like loop such that information about the state of a system can be %%@
sent into the causal past of that system. In classical mechanics this can lead to paradoxes since classical mechanics requires that the basis states %%@
of the system remain fixed within the causality violating region, but quantum mechanics allows a superposition of the basis states which can maintain %%@
a consistent state under the logical operations which would be allowed in a causality violating region. Hence the paradoxes of time travel disappear %%@
when using quantum logic. He also identifies another crucial feature of this arrangement; an observer viewing a companion entering a closed time-like %%@
loop will find that they can enter a state which can be represented as a quantum superposition, but if the observer joins their companion in the %%@
closed time-like loop then they experience one particular outcome with 100\% certainty. This suggests that state vector reduction may be associated %%@
with a change in the reference framework of time.
\par
Using time travel as a way to explain the puzzling features of quantum mechanics has a long history and there have been many attempts to introduce the %%@
notion of information going backwards in
time as a explanation of quantum mechanics, as for example \citet{Aharonov_64,Cramer_jg,Hadley}. The features of quantum computation such as %%@
teleportation and counter-factual computation, ( see \citet{mitcheson_jozsa_2001}), indicate that quantum computation cannot be regarded as just a %%@
large set of classical computations running in parallel. These subtler process rely on the fact that quantum information can appear as if it could be %%@
transmitted backwards in time \citet{penrose_98}.  \citet{Cramer_jg,Cramer_casys} has developed the "Transactional Interpretation of Quantum %%@
Mechanics" using this idea, and \citet*{Antippa_dubois_2004} have shown that \emph{anticipation} in the sense that the future state of a system is one %%@
of the inputs determining the behaviour of the system is a fundamental feature of nature and that such \emph{anticipation} is required for a dynamics %%@
embedded in a discrete spacetime framework.
\par
If a quantum system genuinely could send information backwards in time then we could implement a quantum computation to solve the halting problem, %%@
since we could then create a computation that could send information about its future state \emph{halted} or \emph{not halted} back to an earlier %%@
point in the computation and it's this tantalizing feature of quantum mechanics that has led  \citet{Kieu} and  \citet{Calude_turing_barrier} to %%@
propose that quantum computation could in principle solve the halting problem. However in both cases an infinite number of measurements would be %%@
required. The difficulty is that it's not classical information that goes backwards in time but quantum information and so
it is widely believed that quantum computation does not offer any way to effectively compute this class of
non-computable functions, (see \citet{bennet_and_bernstein} and \citet{Margolus_03}). The issue of the halting of a quantum computer is not resolved, %%@
(see \citet{linden_popescu}) but we should note that Turing's proof \cite{Turing} of the insoluability of the halting
problem relies on the assumption, informed by our experience of classical mechanics, that a system cannot
be in the state of \emph{halted} and \emph{not halted} at the same time. This assumption is not valid in the case of a quantum computer
which can certainly be in a superposition of \emph{halted} and \emph{not halted}. The point that the notion of superposition allows a consistent %%@
answer to problems encountered in closed time-like loops was brought out by \citet{Deutsch_91} and he also argues in %%@
\citet{deutsch_ekert_lupacchini_00} that the physical existence of superposition in quantum systems requires us to extend our understanding of formal %%@
logic to include states which are a superposition of alternatives. 
\par
It is possible here to get bogged down in a lot of discussion about what time is and isn't and what is meant by causal sequence. The usual goal of %%@
such discussions is to find a unifying concept of time which will work for both general relativity and quantum mechanics. We're going to say that such %%@
a thing doesn't exist and take a very liberal notion that time is simply a way to order things so that information in the past is transformed into %%@
information in the present by a computable process. 
In order to make any progress in our understanding we need somehow to introduce non-computability into the scheme.  %%@
\citet{penrose_98,Penrose_ENM,Penrose_SOM} has argued
that the act of quantum measurement introduces some element of non-computability into Nature's functioning. To understand the implications of this we %%@
must make use of computational models of physical law and then find a place in these models for non-computability. So we'll start by examining what we %%@
understand computability to be. First of all the notion of computability was developed in the context of classical mechanics by Alan Turing %%@
\cite{Turing}. The assumption derived from classical mechanics is that if we start with fully specified initial conditions we can progress to other %%@
states in a finite number of discrete steps via a set of rules which can be implemented on some other physical system. This is the basis of the %%@
Universal Turing Machine. We can diagram this as follows;- see figure $(\ref{fig:causal_computation})$ 
\begin{figure}[h]
\centering
\begin{pspicture}(0,0)(8,3)
\rput(0,1.5){\circlenode{A}{\begin{tabular}{c}State \\ \emph{A} \end{tabular}}}
\rput(8,1.5){\circlenode{B}{\begin{tabular}{c}State \\ \emph{B} \end{tabular}}}
\rput(4,2){\psframebox[framearc=0.3]{Some computable process}}
\ncline[nodesep=0.3pt]{->}{A}{B}
\end{pspicture}
\caption{Causal processes assume computability}
\label{fig:causal_computation}
\end{figure}
\par
In this case the only uncertainty comes from insufficient knowledge of the initial conditions and insufficient accuracy during the computable process. %%@
Quantum systems do not follow this scheme since even if we are given enough information about $A$ to specify completely the wave function of the %%@
system at $B$ we still can't predict the outcome of a measurement, only the probability of an outcome. If we add a closed time-like loop then we can %%@
have a system which is not computable in the sense that it can be emulated by a finite process operating according to purely classical physical laws.  %%@
Not all such systems are incapable of being emulated by classical mechanics; only those where the information about the final state of the system that %%@
is sent backwards in time results in a change of the final state of the system such that the output state cannot be represented in a physical system %%@
operating with classical mechanics. We illustrate this in figure (\ref{fig:non_causal_computation})
\begin{figure}[h]
\centering
\begin{pspicture}(0,0)(8.5,3)
\psline{->}(0,0)(7.5,0)
\rput(8,0){\emph{time}}
\rput(0,1.5){\circlenode{A}{\begin{tabular}{c}State \\ \emph{A} \end{tabular}}}
\rput(8,1.5){\circlenode{B}{\begin{tabular}{c}State \\ \emph{B} \end{tabular}}}
\rput(4,2.25){\psframebox[framearc=0.3]{\emph{A} completely determines \emph{B}}}
\rput(4,0.75){\psframebox[framearc=0.3]{\emph{B} completely determines \emph{A}}}
\ncarc{->}{A}{B}
\ncarc{->}{B}{A}
\end{pspicture}
\caption{Non-causal processes are global non-computable but locally computable}
\label{fig:non_causal_computation}
\end{figure}
\par
In the figure (\ref{fig:non_causal_computation}) we've adopted sense of time increasing as we go to the right, and so the connection from \emph{B} to %%@
\emph{A} is identified as information flowing backwards in time. In this case both \emph{A} and \emph{B} determine the output state. But because we %%@
can no longer place \emph{A} and \emph{B} in causal order then we can drop the notion of time completely and simply abstract away this feature of %%@
cyclic dependency. When this feature is encountered in the context of a framework of time then it'll appear is if information can be sent backwards in %%@
time, but there may be other situations where such a framework is absent. Within the context of a framework of time such systems can be said to be %%@
\emph{anticipatory} (see \cite{Antippa_dubois_2004}) in that the future state of a system is part of the input to the system. This can be implemented %%@
as a purely causal process if the final state settles to a single value under the logic implemented by such a causal loop. Indeed this is something we %%@
do all the time when using numerical methods to solve differential equations by taking approximations to some future state and using that as part of %%@
the inputs iteratively to refine the final output. Of course we may end up with unstable conditions such that the algorithm never settles to a single %%@
value and we may choose to model such a situation as a superposition of basis states of the possible solutions. However quantum mechanics does things %%@
a little bit differently, it's not the future state of the system at some arbitrary point that is allowed to become an input but the state of the %%@
system at the future point of measurement, in the measurement basis provided by the observer. In this regard it is useful to adopt the point of view %%@
expressed by  \citet*{Castagnoli_finkelstein_2001} that the algorithm computed in a quantum computation is not fully defined until
we make a measurement. That is, the answer to a computation in the reference frame of the observer is part of the information input into the
computation. This idea is useful because, as a pedagogic exercise, it then becomes quite easy to generate toy
algorithms which embody this property and which then have properties very similar to quantum systems.
\par
To construct these toy models we have to take as an input to the system some feature which is not defined until the answer is computed, and is only %%@
defined in terms of something that an outside observer can measure. This requires that we create some kind of self-referencing statement which can %%@
only be resolved by external observation. What we mean by self-referencing is that the logic of the system refers to itself such that the state of the %%@
system at the point the answer is read out is one of the inputs to the system. If we could create a closed time-like loop then the system could %%@
"observe itself" as if it were an outside observer at some point in the future and use the results of that observation to determine its output. This %%@
would be similar in form to the classically inadmissible case of a Turing machine to solve the halting problem, which could accept a copy of itself as %%@
an input, determine its own resultant state and use that as another input to determine its future state. We could then create a Turing machine that %%@
has no consistent final state according to the laws of classical mechanics. To implement a self-referential process therefore requires a closed %%@
time-like loop, hence a self-referencing process implies non-causality.
\par
For example consider the following self-referential statement, which is defined such that the answer is part
of the question and which can be completed in two possible ways;-
\begin{equation}
\begin{array}{l}
 5\; divided\; by\; the\; number\; of\; characters\; in\; this\; sentence\; plus\; the\; answer\;\\ 
 output\; as\; a\; fraction\; is\;
5 \mathord{\left/
 {\vphantom {5 {}}} \right.
 \kern-\nulldelimiterspace} {} 
 \mathop {\_}\limits_{98} 
 \mathop {\_}\limits_{99} 
 \mathop {\_}\limits_{100}
 \end{array} 
\end{equation}
(The numbers underneath the characters are an index of the characters and are there to allow the reader to
find the two solutions more easily).
\par
Here one of the inputs is something which is only available at the time that the answer is produced, namely the total number of characters in the %%@
sentence. This is a property which has a definite value only in the framework of the observer once the answer has been determined. Obviously this is a %%@
highly contrived example but it does serve a purpose in that it makes much easier to understand the logic. In this instance
there are two possible answers which are consistent with the way the question is formulated.
We can either complete the statement with 99 in which case we get the internally statement
\begin{equation}
\text{\emph{5 divided by the number of characters in this sentence plus the answer
out put as a fraction is 5/99}}
\end{equation}
or with 100, in which case we have a different, but again internally consistent statement
\begin{equation}
\text{\emph{5 divided by the number of characters in this sentence plus the answer
out put as a fraction is 5/100}}
\end{equation}
\par
Having chosen one way to complete the sentence the other possibility immediately becomes invalid, and
from that moment on the only correct answer is the one we selected. The two completed sentences form
the eigenstates of the system and completing the sentence forces the choice of one eigenstate.
This seems like a very trivial instance of a self-referential sentence, and it looks like a simple NP problem
in that there is a two-way propagation of logical implication, and it is in fact an instance of a statisfiability
problem (SAT), where one needs to find a string of bits which can satisfy some logical clause. However it
also has some interesting features in common with a quantum mechanical measurement. The interesting
features are:-
\begin{enumerate}[(a)]
\item %(a)
The "state" of the system is undetermined until the point of "measurement".
\item %(b)
"measurement" constrains the system to output only one of the possible values.
\item %(c)
Once "measurement" has happened the system is consistent with only that one value and the other
value is no longer a possibility.
\item %(d)
In "computing" the output state we note that we don't have all the information required for the
computation until the output has been determined. Therefore the computation cannot be expressed
as a deterministic computation from known inputs and therefore the "output state" will be random, but
evenly distributed between 99 and 100.
\item %(e)
To produce an answer to the statement we must have access to a global property of the system, in this
case the number of characters including the answer, which is information that is locally inaccessible
from within the system.
\item % (f)
The requirement that the answer participates in the formulation of the question severely constrains the
form of the algorithm, the possible values of the starting data and the data at the point of "measurement".
\end{enumerate}
\par
Clearly we can  create a purely causal implementation of something like this as in the case of numerical solutions to differential equations where we %%@
use forward looking approximations to the correct solution and the
computation can settle down to one solution amongst many, but if we change things a little so that we require an input to be something that has to be %%@
determined by an outside observer at the point the answer is read out then we have something that cannot have an implementation by any classically %%@
operating computer. We define a self-referential algorithm as being one which requires one of the inputs to be a property of the system as measured by %%@
an outside observer at the time the answer is determined. With this rule
it's then trivial to create other similar sentences or self-referencing algorithms, for example;-
\begin{equation}
\text{\emph{An input integer value divided by the number of discrete steps required to calculate and output the answer.}}
\end{equation}
\par
Clearly such a number could exist, if there are a finite number of computational steps involved there must
exist at least one rational number that is the input integer divided by the number of steps. And because of
the way we've arranged the logic there may well be more than one. But the system has properties (a)-(f) and
because the logical behaviour of such as system is so close to what we would expect of a quantum measurement
that it is tempting to suggest that a quantum computation could provide a true answer to the question.
\par
There is yet another important feature we can extract from such trivial statements. There are two views
of the computation. One view is from the outside where solutions can be found by exploring all possible
answers to find those that are consistent with the initial statement of the problem. In this case the problem
is simply and NP problem. The other from is the view from inside the computation. This example makes
matters clearer;-
\begin{equation}
\text{\emph{Calculate 5 multiplied by the number of computational steps required to calculate and print out the answer.}}
\end{equation}
\par
Clearly we can create another algorithm that operates outside the computation to do an exhaustive search
of answer states that give a result consistent with the question. In this sense it is an NP-complete problem
and therefore computable. But if we modify the statement so that our external computations are required
to be included in the total number of computational steps as in;-
\begin{equation}
\begin{aligned}
&\text{\emph{Calculate 5 multiplied by the number of computational steps required to calculate and print out the answer }}\\
&\text{ \emph{including any computational steps used by an outside agent.}}
\end{aligned}
\end{equation}
\par
Then such a problem cannot be implemented as a computable algorithm as one of the inputs is information
that is only accessible to an outside observer. So such an algorithm is non-computable.
\par
We have converted an NP-complete problem into a non-computable problem by stepping inside the
self-referential loop and also we can convert a non-computable problem into an NP-complete problem
by stepping outside the loop, where we are in a position to add the information required to compute the
problem. This observation is potentially interesting as \citet*{Castagnoli_Finkelstein_2002_09_084} have argued that
there are non-dynamical forms of quantum computation for which NP=P and Brun \citep{Brun_ctc} has also
shown that a computer that makes use of closed time-like loops can efficiently solve NP problems because
the extra information required can appear as if from nowhere. Once the extra information is available, then
a solution to an NP problem can be verified in P time.

\subsubsection{2.1.1 Is there a level problem here?}
The statement above appears to embody a logic level problem and would therefore be formally illegal. However if it turns out that we need to
introduce logic level problems in order to create toy logical systems that emulate the logic of quantum mechanics then this is surely telling us %%@
something important about the nature of quantum mechanics. 
In fact if we go back to the archetypical logical level problem identified by Russell and Whitehead in \citet{Russel_Whithead}, the statement "Is the %%@
set of all sets that are not members of themselves a member of itself or not?" we can see that the problem arises because classical logic, informed by %%@
classical mechanics, does not allow a superposition of \emph{True} and \emph{False}. But as Deutsch et al \cite{deutsch_ekert_lupacchini_00} have %%@
argued quantum mechanics presents us with physical systems which have states that are not permitted by classical logic so we must extend the %%@
repertoire of permissible logical states to include those which can represent a superposition of \emph{True} and \emph{False}. With such states we can %%@
then present a consistent answer to the question "Is the set of all sets that are not members of themselves a member of itself or not?", the answer %%@
being $\frac{1}{\sqrt{2}}(|YES \rangle+|NO \rangle)$. With such logical states available the reason for the prohibition of logic level errors is %%@
removed. With quantum logic we can allow sets to contain themselves.

\subsubsection{2.1.2 Entanglement as self-referential computation }
\par
However the algorithms we've listed so far seem far removed from any recognizable quantum system. Is
it possible to create a self-referential algorithm that can be identified as a well known quantum system, the
measurement of which would effectively carry out the computation? There is a very well known example
of this, we can set up "entangled statements" obeying the same principle as embodied in the self-referential
statement above, that is, the calculation depends on some feature of the system which is not available until
the point of measurement for example;-
\begin{equation}
\begin{aligned}
X:&\text{\emph{ Calculate a function f(a,y) where a is an input value and y is the result of }}Y\\
Y:&\text{\emph{ Calculate a function f(a,x) where a is an input value and x is the result of }}X\\
\end{aligned}
\end{equation}
We can easily find an equivalent quantum mechanical system such as;-
\begin{equation}
\begin{aligned}
X:&\text{\emph{Set the spin state of particle A be anti-parallel to the result of}} \\ 
  &\text{\emph{measurement of particle B in the measurement basis of Y}}\\
Y:&\text{\emph{Set the spin state of particle B be anti-parallel to the result of}}\\
  &\text{\emph{measurement of particle A in the measurement basis of X.}}\\
\end{aligned}
\end{equation}
\par
This is of course the algorithm that is effectively computed by measurement of spin half entangled particles
in an EPR paradox type experiment, and we know this has a real physical existence. By performing measurement
on a pair of entangled particles we are effectively determining a consistent result for such a self-referencing
computation.
\subsubsection{2.1.3 Extracting the key lesson}
In examining the source of the exponential speed up exhibited by quantum computation \citet{Castagnoli_finkelstein_2001,cast_and_monti_2000} , have %%@
shown that it can be attributed to the non-mechanistic requirement that both the initial
preparation and final measurement are necessary to fully define the problem to be solved. So that the
initial state and the final state play an equal role in determining the parameters of the computation.
Hence the randomness of quantum measurement does not occur by chance and the idea that the measurement 
participates in the problem definition provides a valuable insight to understand the origin of randomness
in quantum measurement. If the final measurement plays an equal role along with initial preparation in
determining the inputs to the problem being solved then it is not possible for the algorithm that is embodied
by the quantum system to be modeled as a purely deterministic evolution from initial starting conditions.
Therefore there can be no deterministic algorithm that can take the initial conditions and produce the
answer, the implication is that each measurement will appear random
and the closest we can get is a statistical calculation of the probability of any one answer. Gregory Chaitin
\cite{Chaitin_limits_book} provides a useful insight from Algorithmic Information Theory.
\begin{quotation}
This is because algorithmic information theory sometimes enables one to measure the information
content of a set of axioms and of a theorem and to deduce that the theorem cannot be obtained
from the axioms because it contains too much information.
\end{quotation}
\par
This is precisely the issue that Castagnoli has identified in quantum computation, that the final measurement
provides the additional information that is not present in the initial conditions that allows us
to fully determine the inputs to the problem. Therefore Castagnoli's insight together with Chaitin's work
on randomness in mathematics allows us to see that the source of randomness in quantum measurement
is the inability of a strictly deterministic computation evolving from fully specified starting conditions to
encompass an algorithm which is self-referencing in the sense that one of the inputs that define the problem is not available
until the answer is determined. The fact that energy has to be provided from an external observer to perform a quantum measurement important because %%@
in the context of computation energy can be identified with information see \citet{Bennet_1982}.
\par As the critical features of quantum systems can be modeled by
self-referencing algorithms and as quantum systems can model self-referential algorithms then quantum
systems very probably are self-referential algorithms. If we are given two black boxes, one which contains
a quantum system the other contains something that can effectively compute a self-referential algorithm
of the type described above, and we are given only the paired values of inputs and outputs for each box,
there is no way we could determine which is which. I'll conjecture that all the apparent paradoxes of
quantum mechanics can be understood as the result of attempting to interpret self-referential logic within
a deterministic framework. 
\par
In this context we can better understand the results of \citet{Kieu} and \citet{Calude_turing_barrier} since they both propose
quantum systems that can solve the halting problem yet each assumes that an infinite number of measurements
can be made on a quantum system. Since a computation to solve the halting problem would require
infinite information then each measurement can be seen as adding the extra information required. However
taking Chaitin's point about being able to determine the information content of axioms and theorems we
can see that there will be certain theorems that are inaccessible by any deterministic computation from any
given finite set of axioms but which will become accessible if we add extra information in the axioms. This implies that
we can define a measure of non-computability for a theorem from a given set of axioms as the number of extra bits
of information that need to be added to the set of axioms to make the theorem computable from that set of axioms.
\par
It should be noted that while the existence of superposition allows a quantum system to represent the
state of a computatation which is formally non-computable according to the rules of classical mechanics, we
do lose that state when we make a measurement. We never perceive superposition directly, it's a notion we
are forced to introduce to make sense of definite measurements. For example if we did create a machine that
could solve the halting problem and then from that created another machine which could self-referentially
examine itself and the halt if it didn't halt, that output state could be implemented in a quantum computer
as a superposition of halted and not halted. But we'd only ever measure it as one or the other. This is why
the complete determination of $\Psi$ would require an infinite number of measurements. This is an important
idea which we will make use of in later on, a quantum system, prior to measurement, can represent the output
state of a computation that is formally non-computable, but physical reality is so constructed that we can
never directly perceive that state. 
\par
We generally take "quantising a system" as meaning to discretize it. Our analysis so far suggests that it means something deeper, "quantising a %%@
system" in fact means to find a mathematical representation of the system that would allow the state of the system at the future point of measurement %%@
to be one of the inputs to the system. This kind of representation can result in a limited number of discrete solutions. \citet*{Antippa_dubois_2004} %%@
have shown that such kind of anticipatory computation is required to model dynamics in a discreet spacetime framework, our analysis suggests that it %%@
may be the other way round, that \emph{anticipation} forces dynamics to be discreet because it adds a stringent contraint on the possible values a %%@
system can take. 
\par
This is the key lesson we want to get from quantum mechanics, the quantum formalism can be used to represent the states of self-referential and hence %%@
non-computable algorithms. This poses a crucial question \emph{"Why does quantum superposition exist and why is physical reality so arranged that we
can never experience it directly?"}
\subsubsection{2.1.4 A signature of self-referentiality?}
We can now identify that the "signature" of self-referentiality, we take self-referentiality of an algorithm to mean that some feature of the system %%@
which is only determined at the point that the answer is read out is an input to the system. This makes such an algorithm formally non-computable,  %%@
and also if we evisage the algorithm as a sequence of events evolving in time it would be non-causal. The essential feature being that that things %%@
only acquire definite values in the context of the answer to a question and cannot have any definite
values except in the context of an answer. 

\subsection{The lesson from general relativity 2.2}
\par
In general relativity non-causal computation is also allowed because general relativity allows the existence
of closed time-like loops. This is the basis of idea presented by \citet{Hadley} that the logic of quantum mechanics
can be derived from general relativity. People often dismiss such solutions to the field equations as being "unphysical", which seems a trifle %%@
reckless since one can always gain useful information from awkward solutions, particularly by asking "why are they unphysical?". General relativity is %%@
a deterministic mathematical model and it shouldn't be possible to model such non-causal computation starting from purely deterministic building %%@
blocks. So there must be something buried in general
relativity that is of inherently non-computable nature. I've suggested that the signature of self-referential computation
is that things acquire values only as part of an answer to a question, general relativity is a theory of space
and time so we should expect that if self-referential logic is buried somewhere within general relativity
then space and time would acquire meaning only as part of the answer to a question. Christian in \citep[chap.~14]{phys_meets_philosophy}
argues that general covariance, carries with it the implication that points in space-time cannot possess any
sense of individuality a priori. As he puts it; -
\begin{quotation}
"Thus, strictly speaking the bare space manifold does not even become 'space-time' with physical
meaning until both the global and local spatio-temporal structures are dynamically determined
along with a metric. Further since spacetime points acquire their individuality in no other way
but as a byproduct of a solution of Einstein's field equations, in general relativity 'here' and 'now'
cannot be part of a physical question, but can only be part of the answer to a question. As Stachel
so aptly puts it \cite{Stachel}, the concepts 'here' and 'now' - and hence the entire notion of local causality -
acquire ontological meaning only \emph{a postereriori}, as a part of the answer to a question."
\end{quotation}
\par
Yet this is precisely the observation that allows us to infer self-referential computation in quantum mechanics, certain aspects of quantum systems %%@
cannot be said to have values until a measurement happens and the determination of the answer to a question participates in the determination
of the problem to be solved. In general relativity it is the observation of an event that participates in the
creation of the whole concept of space and time for the event to occur in. However, there is a difference here in that the notion that future state of %%@
the system becomes an input to the system is absent since we have to define the notion of spacetime first before we know what "future" means so we %%@
have to use the idea of self-reference more abstract way than one which relies on information being sent backwards in time. 
\par
The core problem that frustrates efforts to unite quantum mechanics and general relativity, is that
quantum mechanics does treat space-time as a non-dynamical causal structure in the background that does
allow 'here' and 'now' to be part of the question, points in space-time are assumed to have individuality a
priori, whereas in general relativity they don't, see also \citet{weinstein_2001}. This is the key lesson we take from general relativity which we %%@
need to complement the key lesson we take from quantum mechanics.
\par
\subsection{2.3 Are quantum mechanics and general relativity dual?}
This notion of self-referential logic allows us to identify a potentially interesting duality between general relativity
and quantum mechanics. In quantum mechanics the measurements we make on material bodies carry
the signature of self-referential logic and the background space-time in which the bodies move is strictly
causal, in general relativity space-time itself carries the signature of a self-referential logic and it is the
material bodies in space-time that behave in a strictly causal manner. But such self-referentiality cannot in principle be the the consequence of the %%@
action of any computable process and must therefore be axiomatic. In this case we would have to accept that nature is at root self-referential and yet %%@
quantum mechanics and general relativity introduce a different notion of causal sequence in order to allow a tractable mathematical description of an %%@
underlying
intractable self-referential reality. See figure $(\ref{fig:duality_1})$
\begin{figure}[h]
% This is done using pspicture because in tabular the vertical lines don't connect with the horizontal lines
\centering
\begin{pspicture}(0,0)(12,4)
\psframe(4,3)(8,4)\rput(6,3.5){\emph{General Relativity}}
\psframe(8,3)(12,4)\rput(10,3.5){\emph{Quantum Mechanics}}
% next line
\psframe(0,1.5)(4,3)\rput(2,2.25){\emph{Spacetime framework}}
\psframe(4,1.5)(8,3)\rput(6,2.25){\begin{tabular}{c} Answer is part \\ of the question \end{tabular}}
\psframe(8,1.5)(12,3)\rput(10,2.25){Strictly causal}
% next line
\psframe(0,0)(4,1.5)\rput(2,0.75){\begin{tabular}{c} \emph{Particle interaction in}  \\ \emph{the spacetime framework }\end{tabular}  }
\psframe(4,0)(8,1.5)\rput(6,0.75){Strictly causal}
\psframe(8,0)(12,1.5)\rput(10,0.75){\begin{tabular}{c} Answer is part \\ of the question \end{tabular}}
\end{pspicture}
\caption{Is there a duality between general relativity and quantum mechanics? $\label{fig:duality_1}$}
\end{figure}
\par
Although this is only a hint that such a duality exists, to show it we need to exhibit a function that can transform the concepts embodied in one into %%@
the other, we will assume that it is the case. We know that the differences in the concept of time underlie the problem of finding a theory of quantum %%@
gravity so it would seem a sensible question to ask if these different concepts of time are dual to each other. In the meantime there are two obvious %%@
implications;-
\begin{enumerate}[(1)]
\item
Whatever underlies the operation of quantum mechanics and general
relativity must be purely self-referencing and hence not amenable to being modeled by computable mathematics. So there can be no theory of quantum %%@
gravity and hence no Unified Field Theory. 
\item
All is not lost as we might gain some insight into the distinct forms of those theories by looking at those constraints which must be imposed in order %%@
for each to implement the notion of causal sequence. So now we turn to what those constraints might be.
\end{enumerate}
\section{CAUSALITY PROTECTION AND NON-COMPUTABILITY 3}
\subsection{Non-computability via super-tasks 3.1}
\par
So far we've discussed non-computability in terms of self-referencing systems, that is the answer is part of
the question, and that if information could be sent backwards in time then non-Turing computable action
would be possible. But there exists another way of attaining non-computable action. The precise definition
of a computable action is that it should yield a result in a finite number of steps, if we allow an infinite
number of steps then we can construct a machine that would solve the halting problem. But an infinite
number of steps would take an infinite amount of time so we would not be around to gather the answer.
\citet{Thomson_1954} proposed that one could make a machine that could perform an infinite number of steps in a
finite amount of time. In his example, the Thomson Lamp, if we set up lamp that is switched on for half a
minute, then switched off for a quarter of a minute, then on for an eighth of a minute and so on each switch
over happening exactly after exactly half the period of the previous one then at the end of one minute we
would have switched the light on and off an infinite number of times. But when we look at the light is it on
or off? It cannot be on for that would imply stopping the process after a finite number of steps, likewise it
cannot be off for the same reason. There exists no state in classical mechanics that can describe the state of
the light after an infinite number of switchings. Hence such a thing is impossible. But quantum mechanics
introduces new logical states that allow us to extend logic beyond that which is informed by our experience
of the classical world (see \citet{deutsch_ekert_lupacchini_00}). If we introduce the notion of quantum superposition then the light can be both
on and off at the same time. See figure $(\ref{fig:supertask})$ 
\begin{figure}[h]
\begin{pspicture}(0,0)(8,4)
%\psset{xunit=4cm}
%\psgrid[subgriddiv=1,griddots=10,gridlabels=7pt](0,0)(2,4)
\psset{xunit=1cm}
\psset{linewidth=2pt}
\psline(0,3)(2,3)
\psline(2,1)(3,1)
\psline(3,3)(3.5,3)
\psline(3.5,1)(3.75,1)
\psline(3.75,3)(3.875,3)
\psline(3.875,1)(3.9375,1)
\psline(3.9375,3)(3.984375,3)
\psline(3.984375,1)(3.9921875,1)
\psset{linecolor=gray,linewidth=1pt}
\psline(2,3)(2,1)
\psline(3,1)(3,3)
\psline(3.5,1)(3.5,3)
\psline(3.75,1)(3.75,3)
\psline(3.875,1)(3.875,3)
\psline(3.9375,1)(3.9375,3)
\psline(3.984375,1)(3.984375,3)
\psline(3.9921875,1)(3.9921875,3)
\psset{linecolor=black,linewidth=0.1pt}
%
\rput(1,3){\rnode{A}{\psdots(0,0)}}
\rput(1,4){\ovalnode{B}{ON}}
\nccurve[angleB=270,angleA=90,arrowscale=3]{<-}{A}{B}
%
\rput(2.5,1){\rnode{A1}{\psdots(0,0)}}
\rput(2.5,0){\ovalnode{B1}{OFF}}
\nccurve[angleB=90,angleA=270,arrowscale=3]{<-}{A1}{B1}
%
\rput(4,2){\rnode{A2}{\psdots(0,0)}}
\rput(7,3){\ovalnode{B2}{$\frac{1}{\sqrt{2}}(|ON\rangle+|OFF\rangle)$}}
\nccurve[angleB=180,arrowscale=3]{<-}{A2}{B2}
%
\end{pspicture}
\caption{The end result of a super-task is a superposition of states}
\label{fig:supertask}
\end{figure}
\par
Yet again we have the situation where the quantum formalism is sufficiently rich to be able to model the
logical states that are to be expected as the output of a non-computable task. In this case it's not the
idea of information flowing backwards in time to create causal loops but the notion of a super-task. 
But such a machine cannot be made precisely because of
the physical limits imposed by quantum mechanics on the rapidity of switchings which are possible. In
quantum mechanics the quantity $h$ puts a boundary on how fast we can change the state of something.
This restriction is in addition to the causality protection scheme operated by the finite speed of light. If we
have a system in a superposition of basis states then it looks like the end product of an infinite number
of switchings between the basis states, when we make a measurement we measure one basis state and
the system then looks like only a finite number of changes have occurred. 
\subsection{3.2 Causality protection has two facets}
\par
We hypothesise that although there exists some kind of self-referentiality as fundamental principle of Nature's functioning we do not directly %%@
experience it because there are also causality protection rules which prevent us seeing it.
Causality protection has two facets, $\mathcal{A}$ preventing the use of super-tasks and $\mathcal{B}$ preventing the use of causal loops.
Quantum mechanics and general relativity each have a different notion of time, that is a different notion of how to put things into causal sequence, %%@
so the
obvious question to ask is "Can we recover each of these theories as limit cases by applying these rules over an
assumed non-computable reality?".  
\par
At this point we immediately run into a deep philosophical problem which
I'll highlight and then ignore. An important objection to computational models of physics is that
all such models would appear to assume another underlying physical reality on which to implement the
computation and hence an infinite regression of such physical realities. When we talk about an underlying
non-computable reality we have an even greater problem; such a thing cannot be modeled by any arrangement
of physical particles operating according to normal causal laws. Hence, since our intellectual conceptions are
the products of the operations of material particles operating according to causal laws, we cannot have a
conception of what such a thing might be. Penrose in \citet{Penrose_ENM,Penrose_SOM,penrose_hameroff} has argued that there must be something
of a non-computable nature going on in the brain to produce the experience of consciousness, but until
we have any firm evidence that such a thing could happen we have to conclude that contemplation of
an underlying non-computable reality is outside the bounds our comprehension. This means that any
explanation of how the appearance of causality can arise in a non-computable and hence non-causal reality must, by necessity, be either
incomplete or inconsistent. We can never have a complete theory in the sense we normally understand it,
but must be content getting closer asymptotically to an unattainable goal.
\par
We can however make some general observations about the nature of an non-computable reality which will motivate the ideas to follow;-
\begin{enumerate}[(1)]
\item
A self-referential process could in principle exceed the power of any degree of oracle machine (see \citep[chap. 7.8]{Penrose_SOM}) and
therefore can determine the result of performing any finite computation. Though we should note that the word "process" is not the correct term since %%@
it implies a sequence of things happening in time. 
\item
Since such a process can determine the truth or falsehood of any theorem then it is complete, and so by 
Gödel's First Incompleteness Theorem \cite{Godel} it must be inconsistent.
\item
Being inconsistent theorems $P$ and \emph{not} $P$ are both provable by the system.
\item
Hence the system can negate itself.
\item
The negation of self-reference is non self-reference, so causal sequences are produced from the system, and because all theorems can be checked by the
process, all possible causal sequences are outputs. 
\end{enumerate}
\par If we start with an abstract notion of self-referentiality then we can by some means evolve all possible causal sequences. We now make the
simplifying assumption that reality is inherently self-referential, and thus not capable of being represented as a causal sequence, but we want to %%@
find an approximate representation of it as if it were a causal sequence. 
\par
The first approximation is that we're not going to be rigorously specific about the possible states of such a non-computable computation, we'll not %%@
specify what we mean by states or states of what. We will just assume that there are patterns
of information, and that some patterns can be transformed into others by a computable, and hence causal,
process and others cannot. From any axiom system certain theorems are provable and some are not unless
we extend the axiom system. This is the case that \citet{Chaitin_Randomness_in_Arithmetic} has pointed out. So without being more
specific, we'll just assume that the output of a non-computable computation contains states that can be
put into causal order and states that can't. If we attempt to model an underlying non-causal reality with
either of the causality protection rules in place we have to find a place in that model for states that will
arise if the rule were violated. 
\par
If the idea that time, or causal sequence, is something superimposed over
an underlying non-causal reality is to be any good then as a bare minimum it should give us some insight
into the nature of space and perhaps even an understanding of non-locality. So with these caveats in mind
we look at the consequences of each rule.
\subsection{3.3 Preventing Super-Tasks}
\par
We'll define our first rule as;-
\begin{equation}
\text{Rule }\mathcal{A} \text{ we can't we can't access the output of a super-task}
\end{equation}
Since a super-task goes through an infinite number of steps in a finite amount of time measured by reference to an external observer, then we prevent %%@
a super-task from happening by enforcing a rule that there is a finite maximum number of
configurations that a system can go through in a finite period of time measured according to the external reference
frame. Following \citet{Margolus_03} we define $\Omega(t)$ the number of configuration changes in a unit of time.
This is a dimensionless number, but if we call this number action we can then create another concept energy
which is the maximum rate at which configurations can change. We then have a constant of proportionality,
let's call it $K$ which puts a boundary on the number of configuration changes in a unit time via
\begin{equation}
Et=K\Omega(t)
\end{equation}
\par
So that $K$ has the dimension of $energy \times time$ since $\Omega(t)$ is dimensionless. The actual equation in our world 
is
\begin{equation}
Et=\frac{h}{2}\Omega(t)
\end{equation}
So that $K=h/2$, hence we can see that the operation of rule $\mathcal{A}$ demands a fundamental constant which has the units of action. But it also %%@
requires us to have an external reference framework of time.
\par
In this instance we are creating a representation of the world by putting this
rule in place over an underlying self-referential, hence non-computable reality so we have to have a way to
represent the states of non-computable processes. A non-computable process could decide the result of a formally undecidable proposition, so we have %%@
to have a way to represent a formally undecidable proposition. Under this scheme we represent a non-computable process as
a violation of rule $\mathcal{A}$, if we violate rule $\mathcal{A}$ then we could have a super-task so in this case we represent a formally %%@
undecidable proposition as a superposition of possible outcomes, like the state
of Thomson's Lamp at the end of one minute. With the rule $\mathcal{S}$ in place we interpret an underlying non-computable
reality as if it we a causal sequence of events evolving in a background reference framework
of time and a constant which has the dimension of action which serves to prevent the output of a super-task ever being accessed, we represent the %%@
violation of the rule as a superposition of states. See figure $\ref{fig:qm_1}$
\begin{figure}[h]
\begin{pspicture}(-0.5,-0.5)(6,6)
% qm picture
\rput(2,5){\emph{Quantum Mechanics}}
\psline{->}(0,0)(0,4) 
\rput(-0.85,3){\rnode{B1}{$\frac{1}{\sqrt{2}}|FALSE\rangle$}}
\psline{->}(0,0)(4,0)
\rput(3,-0.3){\rnode{B2}{$\frac{1}{\sqrt{2}}|TRUE\rangle$}}
%\nccurve[angleB=180,arrowscale=2]{<-}{A2}{B2}
\psline{->}(0,0)(3,3)
\rput(4,4){\ovalnode{B}{$\frac{1}{\sqrt{2}}(|TRUE\rangle+|FALSE\rangle)$}}
\rput(3,3){\rnode{A}{\psdots(0,0)}}
\psline[linestyle=dotted](0,3)(3,3)
\psline[linestyle=dotted](3,0)(3,3)
\nccurve[angleB=270,arrowscale=2,nodesep=0.1]{<-}{A}{B}
\end{pspicture}
\caption{Quantum mechanics represents an undecidable proposition as a super-position \newline of alternative outcomes }
\label{fig:qm_1}
\end{figure}
\par
Hence applying this rule allows us to start recovering quantum mechanics from the notion of self-referentiality.
It also reveals an interesting conclusion, which is not anticipated in the quantum formalism; simply by
stating rule $\mathcal{A}$ as "there will be no super-tasks" we deduce that at any point in time there can only be a finite
number of prior configuration changes. Therefore interpreting a non-computable reality using this rule
makes it look like causal processes began at a finite time in the past.
\subsection{3.4 Preventing causal loops}
\par
We define our second rule as;-
\begin{equation}
\text{Rule }\mathcal{B} \text{ we can't we can't access the output of a causal loop }
\end{equation}
In this case we enforce a rule over an assumed non-computable reality such that if rule $\mathcal{B}$ could be violated we could create a causal loop. %%@
This rule does not require a global reference framework of time; it simply requires that locally we only ever experience time as a strict causal %%@
sequence in the changes of the configuration of things. This has interesting consequences, in special and general relativity if we could send signals %%@
at superluminal speed then we can create a causal loop. See figure $(\ref{fig:ftl})$.
\begin{figure}[h]
\centering
\begin{pspicture}(-2,0)(6,3)
\psset{arrowscale=3}
\psline{->}(3,0)(1.5,1.5)
\psline{->}(1.5,1.5)(0,3)
\psline{->}(0.5,2.5)(4.5,1.5)
\psline{<-}(2.5,0.5)(4.5,1.5)
\psline{->}(3,0)(6,3)
\rput(0.5,2.5){\psdots[dotscale=2](0,0)}
\rput(0,2.5){A}
\rput(4.5,1.5){\psdots[dotscale=2](0,0)}
\rput(4.8,1.5){B}
\rput(-1.75,1.75){\begin{tabular}{c}A sends a superluminal\\ signal to B \end{tabular}}
\rput(7,1){\begin{tabular}{c} B responds with \\ a superluminal  signal \\ to the past of A \end{tabular}}
\end{pspicture}
\caption{Superluminal signaling can be used to create a causal loop}
\label{fig:ftl}
\end{figure}
\par
But superluminal signaling means that we can communicate between space-like separated events, so in this scheme we represent a formal undecidable %%@
proposition as an event outside the light cone, see figure $(\ref{fig:duality_2})$.
\begin{figure}[h]
\centering
\begin{pspicture}(0,0)(6,6)
%gr picture
\rput(2,5){\emph{Einsteinian Causality}}
\psline{->}(0,0)(4,4)\psline{->}(4,0)(0,4)
\psline{->}(0,2)(0,3)\psline{->}(0,2)(1,2)
\rput(0.4,2.8){\emph{time}}
\rput(0.5,1.8){\emph{space}}
\rput(2.25,2){\rnode{A}{\psdots(0,0)}}
\rput(5,3){\rnode{B}{\begin{tabular}{c}\emph{Undecidable}\\ \emph{Proposition} \end{tabular}}}
\nccurve[angleB=180,arrowscale=2,nodesep=0.1]{<-}{A}{B}
\end{pspicture}
\caption{In special and general relativity a undecidable proposition is represented as a space-like separated event}
\label{fig:duality_2}
\end{figure}
\par
In this picture the notion of space-like separation is born from the need to represent non-causality as if we could create a closed time-like loop.
Straight away by putting figure $(\ref{fig:qm_1})$ alongside figure $(\ref{fig:duality_2})$ we can see why spatial separation in quantum mechanics %%@
does not carry the same meaning of causal separation that it does in special and general relativity. Under rule $\mathcal{A}$ causality protection is %%@
ensured by the quantity $h$ under rule $\mathcal{B}$ causality protection is ensured by the quantity $c$, the speed of light. Therefore the %%@
representation of reality we get from $\mathcal{A}$ does not see $c$ as a limit on the speed of information transfer between space-like separated %%@
measurements. Similarly the representation of reality under rule $\mathcal{B}$ does not use $h$ and so general relativity is formulated as if $h$ does %%@
not exist.
\par
By looking at figure $(\ref{fig:ftl})$ we can see something else, the arrows have to be lined up head to tail. Being able to line all the arrows up %%@
head to tail like this means that we can parallel transport our vectors around the loop. The property of a circle that allows us to do parallel vector %%@
transport is called "parallelisability" and the circle is a one dimensional parallelisable sphere. This sphere is called $S^1$, we get the unit sphere %%@
$S^1$ by taking  exponentials of unit modulus complex numbers. We now observe that ;-
\begin{enumerate}[(1)]
\item
If we count the number of configuration changes required in a computation that proceeds from a set of axioms to a given theorem as a real quantity %%@
then we can use that as a measure of time. 
\item
We can measure the degree of non-computability of a theorem from a given set of axioms; it's the amount of information we need to add to an axiom %%@
system in order to be able to derive the theorem by computable means. The measure of space-like distance between space-like separated events is the %%@
amount of information that has to be added to allow one event to become causally connected to another.
\item
A causal loop can add information as if from nowhere.
\item
We can generate a loop by taking the exponential of a complex number then the amount of information that needs to be added can be represented by a %%@
complex number. 
\item
Then we can we measure space-like distances using complex numbers so that with one spatial dimension the metric has signature $(+-)$.
\par
Now that is not a proof, the only justification for it at present is that it's how Nature does things. This is a major open problem in this line of %%@
argument. We need to prove that if we represent causal steps in a computation with real numbers we can represent non-causal steps using complex %%@
numbers. But since that is how Nature does things we will assume that such a proof will be available at some point and move on to the next step in the %%@
argument. 
\item
There are two other continuous spheres that have the parallelisability property, they are $S^3$ and $S^7$. Each of the continuous parallelisable %%@
spheres is associated with a complex normed division algebra.  So that the unit circle $S^1 \equiv e^{z}\text{, } \{z \in \mathbb{C}:|z|=1 \}$ and %%@
similarly with the quaternions and octonions to get $S^3$ and $S^7$, (see \citet{baez_2002} for a comprehensive introduction to the octonions). This %%@
means that we also have to exclude the possibility of creating causal loops on $S^3$ and $S^7$. Which in turn means we now have to represent causal %%@
loops by parallel vector transport on these spheres. But these spheres are the equivalent to the exponentiation of the quaternions and the octonions. %%@
The quaternions are defined by
\begin{equation}
i^2=j^2=k^2=ijk=-1
\end{equation} 
from which we deduce that
\begin{equation}
ij=k
\end{equation}
so we only need two base quaterions to create the full set. 
\item
Similarly with the octonions we need only three base octonions to generate all of them and this means that we can generate any rotation on $S^7$ from %%@
three base rotations. Hence the information required to create a causal loop by parallel transporting our vectors around $S^7$ is contained in just %%@
three quantities.
\item
If we measure computable steps using real numbers and we use complex numbers to count non-computable steps then since rotations on the largest sphere %%@
can be represented by the exponentiation of just three independent complex numbers then we must measure space-like separation using at most three %%@
dimensions. Furthermore since we are counting non-computable steps using complex numbers these dimensions must be complex with respect to the time %%@
dimension so that the metric has signature $(+---)$. 
\end{enumerate}
\par
In summary;If we could send information between space like separated events then we could create a causal loop. A causal loop requires the %%@
parallelisability property. Rotations on the largest parallelisable sphere can be specified with just three independent pieces of information and %%@
therefore space has three dimensions.
\par
We also see that we've had to make use of the notion of computability and computable sequences in order to define our rule $\mathcal{B}$. From %%@
\citet{Margolus_03} we can understand computation, i.e. changes in the configuration of things, in terms of energy and momentum and so to define a %%@
metric on the notion of space and time which springs into existence under the operation of rule $\mathcal{B}$ we have to define it in terms of energy %%@
and momentum. This brings us very close to Einstein's field equations, although not quite there. We can see that we must have three spatial %%@
dimensions, with one of time and a metric with signature $(+---)$ and that it's going to involve defining space-time in computational terms, i.e. in %%@
terms of energy and momentum, but how we get the metric isn't clear. This is work in progress.
\par 
That we are able to find a pointer to a plausible explanation for the observed number of dimensions of space and the non-locality of quantum phenomena %%@
with such little effort suggests that the view that time is a notion of causal sequence super-imposed on a non-causal reality using a causality %%@
protection rule may be good one to pursue further. However so far we have only a plausible explanation, not a rigorous one, but it is none the less %%@
encouraging. It shows that the dimensionality of space is not something determined by the flavour of the month in particle theories but is completely %%@
determined by the actual meaning of space-like separation. Space like separation has the hidden implication that if we could send signals between %%@
space-like separated events then we could create causal loops and it is this that determines the number of space dimensions. If we have a theory that %%@
requires more than three apparent spatial dimensions then these extra dimensions cannot have the meaning of causal separation that they do in the %%@
relativistic viewpoint.
\par
This insight reveals something quite unexpected, if we translate it into the framework of Algorithmic Information Theory it implies that the %%@
information content of new axioms added to an axiom system can be measured in at most three independent ways and also that we can measure the degree %%@
of non-computability of a theorem from a given finite set of axioms by counting in imaginary numbers. Proving this is a key open problem. 
\par
Bearing in mind that much work remains to be done I'll conjecture that the relationships between quantum mechanics, general relativity and a %%@
fundamental purely self-referential and hence non-computable reality can be summarised in figure $(\ref{fig:qm_gr_1})$
\begin{figure}[h]
\begin{pspicture}(0,0)(6,4)
\rput(3,3.5){\rnode{A}{\psframebox[framearc=0.3]{Underlying Self-Referential Reality}}}
\rput(1,0.5){\rnode{B}{\psframebox[framearc=0.3]{Quantum Mechanics}}}
\rput(5,0.5){\rnode{C}{\psframebox[framearc=0.3]{General Relativity}}}
\ncline[nodesep=0.2]{->}{A}{B}
\ncline[nodesep=0.2]{->}{A}{C}
\rput(1.75,2){\psframebox*{\begin{tabular}{c}Enforce \\ rule  $\mathcal{A}$ \end{tabular}}}
\rput(4.25,2){\psframebox*{\begin{tabular}{c}Enforce \\ rule  $\mathcal{B}$ \end{tabular}}}
\end{pspicture}
\caption{Conjectured relationship between quantum mechanics \newline general relativity and an underlying reality }
\label{fig:qm_gr_1}
\end{figure} 
\subsection{3.5 Further development of the idea}  
We'll now step on to more tentative areas. If we look at the implementation of rule $\mathcal{A}$ we see we have to define it with reference to an %%@
external global framework of time, but an external framework of space and time seems to require the implementation of rule $\mathcal{B}$. Yet to %%@
define rule $\mathcal{B}$ we have to assume that computations can only proceed in finite steps so our rules are better defined as;-
\begin{equation}
\text{RULE:}\mathcal{A}\text{ -  Assuming rule }\mathcal{B}\text{ is in place then in addition we can't access the output of super-tasks}
\end{equation}
and
\begin{equation}
\text{RULE:}\mathcal{B}\text{ -  Assuming rule }\mathcal{A}\text{ is in place then in addition we can't access the output of causal loops}
\end{equation}
\par
This then tells us what this "underlying self-referential reality" actually is, it's the totality of the physics generated by general relativity and %%@
quantum mechanics together. Expressing it in the form of two separate concepts each of which is purely causal in its own domain allows us to use %%@
standard computable mathematics. When we put them together we get a system where the concepts underpinning each one depend on the other. The notion of %%@
there being finite number of states  a system can go through in a finite amount of time, which is essential for the creation of the spacetime %%@
framework of general relativity requires quantum theory to ensure that only a finite number of states are possible. Quantum theory has to make use of %%@
the framework of space and time that is ultimately derived as a limit case from the relativistic view. Hence if we put them together we get a fully %%@
self-referential system which is mathematically intractable. The case is closely analogous to the case shown in figure %%@
$(\ref{fig:non_causal_computation})$ where each segment of a loop in isolation is computable though the whole loop is not. The two rules together %%@
therefore represent a more detailed view of this "self-referential reality" which we can represent in this diagram $(\ref{fig:first_expansion})$
\begin{figure}[h]
\begin{pspicture}(0,0)(10,6)
\cnode(1.5,3){1.5}{A}
\cnode(8,3){2}{B}
\ncline[arrowscale=3]{->}{A}{B}
\rput(4.5,3.25){Is detailed as}
\rput(1.5,3){\begin{tabular}{c}Pure \\ Self-referentiality \end{tabular}}
\rput(8,3){
\rput(0,1){\ovalnode{A1}{Rule $\mathcal{A}$}}
\rput(0,-1){\ovalnode{B1}{Rule $\mathcal{B}$}}
\nccurve[arrowscale=3]{->}{A1}{B1}
\nccurve[angleB=180,angleA=180,arrowscale=3]{->}{B1}{A1}
}
\end{pspicture}
\caption{Rules $\mathcal{A}$ and $\mathcal{B}$ together form a more detailed version of the self-referential system}
\label{fig:first_expansion}
\end{figure}
\par
The consequences of this arrangement are;-
\begin{enumerate}[(1)]
\item Since $\mathcal{A}$ and $\mathcal{B}$ form a self-referencing logical loop, and one underlies quantum mechanics and the other general relativity %%@
then a theory of quantum gravity would have to be fully self-referential which means it could model a mechanism to solved the halting problem which %%@
isn't possible, therefore there is no theory of quantum gravity.\\
\item In cases where we have infinite energies then we can't assume $\mathcal{A}$ so rule $\mathcal{B}$ breaks down. Since rule $\mathcal{B}$ is %%@
needed to construct the notion of space-like separation then the notion of space breaks down. This is the point at which general relativity breaks %%@
down.
\item At the point of measurement of a quantum system, since the system is defined with the value at measurement as part of the inputs then we have an %%@
apparent causal loop, which means that rule $\mathcal{B}$ is invalid and so rule $\mathcal{A}$ breaks down, and since rule $\mathcal{A}$ is required %%@
to create the notion of super-position then we lose super-position of possibilities.
\item
Rule $\mathcal{A}$ requires a global notion of time, while rule $\mathcal{B}$ only a local one, so as we go round the loop some things will be %%@
invariant under global-local transformations hence the laws of nature that arise from the interaction of these rules must be expressed as gauge %%@
theories.
\item
\begin{enumerate}[(a)]
\item
In the relativistic case the number of degrees of freedom in space is constrained to three by the operation of rule $\mathcal{B}$, but in the quantum %%@
case space-like separation is not constrained to three dimensions. We can have an infinite number of degrees of freedom and only some of these can be %%@
mapped into the framework of space used when we make a measurement. How are the others to be represented? 
\item
In terms of the observer framework derived from rule $\mathcal{B}$ we can interpret these extra degrees of freedom as extra spatial dimensions with %%@
the peculiar property that $c=\infty$ which makes them look like tiny causal loops which we can represent as a fibration over Minkowski space using %%@
the parallelisable spheres. 
\item
But these spheres arise form the exponentiation of the hypercomplex division algebras and a result due to \citet{dixon_G_94,dixon_99} and other work %%@
by \citet{smith_t_95} shows that we can derive the $U(1)\times SU(2) \times SU(3)$ symmetry of the Standard Model from the hypercomplex division %%@
algebras. 
\item
Since the interaction between rules $\mathcal{A}$ and $\mathcal{B}$ when we make a measurement requires the transformation from a global concept of %%@
time to a local concept of time then we should expect that our measurements can be interpreted in the form of gauge transformations with $U(1)\times %%@
SU(2) \times SU(3)$ symmetry. So there is a case to be made that the symmetry of the Standard Model is the inevitable consequence of observing a %%@
self-referential totality as if it were a locally a causal sequence. 
\par
Although the work on the relationships between the symmetry groups of the Standard Model and the division algebras was done a long time ago it has %%@
languished unregarded because it's hard to fit it into research programs attempting to develop a Unified Field Theory, or even a Grand Unified Theory. %%@
However the arguments presented here suggest that such programs may well fail and that the significance of parallelisable spheres is greater than %%@
previously thought. 
\end{enumerate}
\end{enumerate}
\par
But such a derivation would have peculiar consequences;-
\begin{enumerate}[(a)]
\item
In this a picture the symmetry group ultimately derives from rotations of $S^7$ which don't form a group because octonion multiplication is %%@
non-associative. The symmetry group springs into existence fully formed rather than via symmetry breaking from some higher order group such as %%@
$SU(5)$. The implication is that there is no Grand Unified Theory, which fits the current data, and can be refuted by experiments due to be performed %%@
at CERN.
\item
If the Standard Model symmetry group is the consequence of the interaction between causality protection rules, then could particle mass be another %%@
consequence? A natural derivation of particle mass could come from finding a way to define the metric of general relativity in computational terms. %%@
Because entanglement is a computational resource in quantum systems it could be that the capacity for particles to engage in entangled ensembles is %%@
what gives rise to mass. The idea that there may be a relationship between entanglement and mass has been suggested by \citet{Vedral_2004}. However he %%@
proposes that the entanglement could be between virtual Higgs bosons and particles, though of course entanglement with any other virtual particle will %%@
do just as well and then there would be no need for a Higgs boson. In our suggested scheme, a measure of space like distance between events is the %%@
amount of information that has to be added to allow one event to become causally connected to another. If there are particles present which are %%@
engaging in entanglements with each other and the vacuum then this will obviously affect the amount of information separating events and hence affect %%@
the metric.   
\end{enumerate}
\section{4 CONCLUSIONS}
\par
In this essay I've explored some of the issues raised by the question "Why is the quantum formalism so uniquely adapted to representing %%@
self-referential, non-computable states?". We've seen that this issue has been addressed already in many different guises. By taking the lesson %%@
seriously and assuming that it is telling us something important about the nature of reality we've been able to formulate arguments that may lead to %%@
solutions of otherwise intractable problems. 
\par
I've had to jump logical chasms with arguments that are are not properly justified and have left open problems, but the broad outlines of a theory %%@
that starts to emerge from the mist are nevertheless very interesting; it begins to match what we see in the world and seems to give some insight into %%@
hard problems. The idea that
self-referentiality may be a foundational principle is not novel, it has for example been been mooted by \citet{hawkin_2003}
in despair at ever finding a mathematically consistent Unified Field Theory. A similar idea has been presented in  a different form by
\citet{Barbour_EOT}. Barbour argues that since the notion of time is so different in quantum mechanics and
general relativity that we should abandon the notion of time altogether. Since this would imply that causal
sequence is not a fundamental property of nature this seems to suggest that non-causality is fundamental
instead. 
\par
The early attempts to create a theory of quantum gravity ran into problems with causality and infinite energies. Subsequent programs in this area have %%@
attempted to create theories which avoid these problems and at the same time match the world we live in. These programs have not been as succesful as %%@
people have hoped (see \citet{smolin_03}). The suggestion that self-referentiality may be attained by causality violation or infinite energy implies %%@
that if it is to be regarded as a foundational principle then the problems encountered in the early work were actually the right answers.   
\par
However as we have seen the case is not a bad as it looks, for while it seems that a mathematically consistent
Unified Field Theory cannot exist, we can at least get some understanding of why
our limited theories have the forms that they do. Because getting a consistent result from a self-referential algorithm places very strong constraints %%@
on the possible solutions, the requirement that the sum total of physical laws form a self-referential system places very strong constraints on the %%@
form those physical laws can take if we want the physical theory for each segment of the logical loop to present its own consistent view of Nature. %%@
The arguments gathered so far point to there being profound reasons for the world being the way it is, not as some have suggested merely a lucky %%@
selection from maybe $10^{500}$ different string theories (see \citet{susskind_2003,susskind_2003_ns}). It may be that if we pursue this idea we will %%@
find that the free parameters of the Standard Model have to take the values that they do in order to preserve self-reference in the totality of %%@
physical laws.  An added bonus is that since a self-referential system can generate information as if from nowhere then it can also generate the %%@
physics required to implement itself from nowhere. So, seemingly out of nowhere, we would get a set of physical laws that are individually purely %%@
causal, but collectively non-causal. When viewed from a narrow perspectives of limited theories the physics would be purely causal with three %%@
dimensions of space, a spacetime metric defined in terms of computations, quantum non-locality and a gauge theory using a $U(1) \times SU(2) \times %%@
SU(3)$ group. But if we attempt to create a single mathematical model of all of physical law together it would contain causality violations and %%@
infinite energies. There remain major technical problems and great philosophical problems in this approach. The philosophical problems I propose to %%@
ignore and I'll just list the major technical problems to be resolved. 
\subsection{4.1 Open problems}
\begin{enumerate}[(1)]
\item % (1)
The idea that if we count computable steps with numbers on the real axis, we must use complex
numbers to count non-computable steps is crucial. This needs to be placed on a rigorous footing. The suggestion, hinted at from Nature and supported %%@
with knowledge of the parallelisable spheres, that we must measure the degree of non-computability in at most three linearly independent ways is very %%@
odd, entirely unexpected, and needs more study. 
\item % (2)
A super-task and a causal loop can both compute what is otherwise non-computable, we need to show
that they are in fact equivalent in this sense. I've simply stated without proof that we can represent the end state of a super-task as a %%@
superposition of outcomes. Proof is required. 
\item % (3)
I've suggested that there is a duality between general relativity and quantum mechanics, we need to exhibit a transformation that maps one to the %%@
other.
\end{enumerate}
\bibliographystyle{plainnat}
\bibliography{quantm_self_reference_3}
\end{document}

